\paragrafos{Θέσεις ευθειών στο επίπεδο: τεμνόμενες και παράλληλες}{Β.1.9, «Θέσεις ευθειών στο επίπεδο»}

\begin{anagke}
Ο παππούς δείχνει πώς θα φυτευτούν οι σειρές με τα καρότα στον κήπο: «Κάθε σειρά θα είναι ευθεία και όλες θα έχουν την ίδια κατεύθυνση. Καμιά δεν θα ακουμπά την άλλη όσο κι αν τις συνεχίσετε». Η Αγγελική ξαφνικά σκέφτεται τις γραμμές στο σχολικό γήπεδο, τις γραμμές στο τετράδιο και τις ράγες του τρένου: «Δεν συναντιούνται ποτέ».

Ο Ηλίας δείχνει δύο μονοπάτια του κήπου που διασταυρώνονται: «Αυτά ακουμπούν σε ένα σημείο». «Μα αν συνεχίσεις ένα μονοπάτι που δεν διασταυρώνεται, μπορεί να διασταυρωθεί αργότερα!» διαμαρτύρεται ο Στέλιος. Η Πολυξένη προτείνει ένα τεστ: «Σκέψου τις ευθείες απεριόριστες. Δύο ευθείες του ίδιου επιπέδου είτε έχουν ένα κοινό σημείο είτε δεν έχουν κανένα».

Στην παράγραφο αυτή ονομάζουμε τις δύο δυνατές θέσεις: τεμνόμενες και παράλληλες. Και ο Ευκλείδης θα μας δώσει ένα μοναδικό αξίωμα για το πώς φέρνουμε παράλληλη από ένα σημείο.
\end{anagke}

\begin{lexi}{παράλληλες, τεμνόμενες}{παράλληλες, τεμνόμενες}
\lexiword{Παράλληλες} είναι οι ευθείες που πηγαίνουν «δίπλα-δίπλα» (παρ' άλληλες), «η μία πλάι στην άλλη», χωρίς ποτέ να συναντιούνται. \lexiword{Τεμνόμενες} είναι οι ευθείες που «τέμνονται», δηλαδή σταυρώνουν.
\end{lexi}

\begin{metafrasi}
Δύο ευθείες χωρίς κοινό σημείο. & παράλληλες: $\varepsilon_1\parallel\varepsilon_2$\\
Δύο ευθείες με ένα κοινό σημείο. & τεμνόμενες\\
Το κοινό τους σημείο. & σημείο τομής\\
Τμήματα πάνω σε παράλληλες ευθείες. & $\mathrm{ΑΒ}\parallel\Gamma\Delta$\\
\end{metafrasi}

\tmima{Η θεωρία}

\begin{orismos}{παράλληλες και τεμνόμενες ευθείες}
Δύο ευθείες του ίδιου επιπέδου λέγονται \textbf{παράλληλες} αν δεν έχουν κανένα κοινό σημείο, όσο κι αν προεκταθούν. Γράφουμε $\varepsilon_1\parallel\varepsilon_2$. Λέγονται \textbf{τεμνόμενες} αν έχουν ένα (και μόνο ένα) κοινό σημείο, που λέγεται \textbf{σημείο τομής}.
\end{orismos}

Δύο ευθείες του ίδιου επιπέδου ή είναι παράλληλες ή τέμνονται (αν έχουν δύο κοινά σημεία ταυτίζονται, αφού από δύο σημεία περνά μία ευθεία). Δύο τμήματα λέγονται \textbf{παράλληλα} όταν βρίσκονται σε παράλληλες ευθείες.

\begin{center}
\begin{tikzpicture}[scale=0.9]
\draw[thick,{Stealth}-{Stealth}] (0,0)--(4,0) node[right] {$\varepsilon_1$}; \draw[thick,{Stealth}-{Stealth}] (0,1.3)--(4,1.3) node[right] {$\varepsilon_2$};
\node at (2,-0.6) {$\varepsilon_1\parallel\varepsilon_2$};
\begin{scope}[xshift=7cm]
\draw[thick,{Stealth}-{Stealth}] (0,0)--(4,1.6); \draw[thick,{Stealth}-{Stealth}] (0,1.5)--(4,-0.1); \fill (2,0.77) circle (2pt);
\node at (2,-0.6) {τεμνόμενες};
\end{scope}
\end{tikzpicture}
\end{center}

\begin{idiotita}{αξίωμα παραλλήλων (Ευκλείδη)}
Από ένα σημείο $\mathrm{Α}$ εκτός ευθείας $(\varepsilon)$ διέρχεται \textbf{μία και μοναδική} ευθεία παράλληλη προς την $(\varepsilon)$.
\end{idiotita}

Το αξίωμα αυτό ανήκει στο έργο του Ευκλείδη, τα «Στοιχεία», που γράφτηκε στην Αλεξάνδρεια περίπου το $300$ π.Χ. Για πολλούς αιώνες οι μαθηματικοί προσπαθούσαν να το αποδείξουν από τα άλλα, και τελικά κατάλαβαν ότι είναι αναγκαίο: υπάρχουν και γεωμετρίες όπου δεν ισχύει.

Μπορούμε να φέρουμε παράλληλη με \textbf{γνώμονα και χάρακα}: ακουμπάμε τον χάρακα στη δοσμένη ευθεία $(\varepsilon)$ και τον γνώμονα πάνω στον χάρακα. Ολισθαίνουμε τον γνώμονα κατά μήκος του χάρακα μέχρι να περάσει από το σημείο $\mathrm{Α}$ και χαράσσουμε. Η ευθεία που προκύπτει είναι παράλληλη.

\begin{prosoxi}
\begin{itemize}[leftmargin=1.2em,itemsep=0.3ex,topsep=2pt]
\item Για να είναι παράλληλες δύο ευθείες δεν αρκεί να μην συναντιούνται στο σχέδιο: πρέπει να μην έχουν κοινό σημείο όσο κι αν προεκταθούν.
\item Δύο παράλληλα τμήματα δεν χρειάζεται να έχουν το ίδιο μήκος.
\item Στο επίπεδο δεν υπάρχουν δύο ευθείες που να μην είναι ούτε τεμνόμενες ούτε παράλληλες (στον χώρο υπάρχουν, οι ασύμβατες).
\item Αν μια ευθεία είναι παράλληλη σε δύο άλλες, αυτές είναι παράλληλες μεταξύ τους.
\end{itemize}
\end{prosoxi}

\tmima{Λυμένα παραδείγματα}

\paradeigma{Παράλληλες ή τεμνόμενες;}
\tixero Οι ευθείες: τα όρια του γηπέδου, οι ράγες, δύο μονοπάτια που διασταυρώνονται.\par
\tithelo Να πω ποιες είναι παράλληλες.\par
\possindeo Παράλληλες: οι ράγες και τα απέναντι όρια του γηπέδου. Τεμνόμενες: τα μονοπάτια.\par
\apantisi Ράγες και απέναντι πλευρές: παράλληλες. Μονοπάτια: τεμνόμενες.\par
\elegxos Οι παράλληλες δεν έχουν κοινό σημείο.

\paradeigma{Φέρνω παράλληλη}
\tixero Μια ευθεία $(\varepsilon)$ και ένα σημείο $\mathrm{Α}$ εκτός της.\par
\tithelo Να φέρω την παράλληλη.\par
\possindeo (1) Τοποθετώ τον χάρακα πάνω στην $(\varepsilon)$. (2) Τοποθετώ τον γνώμονα πάνω στον χάρακα, με μια πλευρά κατά μήκος του. (3) Σύρω τον γνώμονα μέχρι η άλλη του πλευρά να περάσει από το $\mathrm{Α}$. (4) Χαράσσω.\par
\apantisi Η ευθεία από το $\mathrm{Α}$ είναι η μοναδική παράλληλη προς την $(\varepsilon)$.\par
\elegxos Μετρώ δύο φορές την απόσταση των ευθειών: βγαίνει ίδια.

\paradeigma{Πόσα σημεία τομής;}
\tixero Τρεις διαφορετικές ευθείες του ίδιου επιπέδου.\par
\tithelo Πόσα το πολύ και πόσα το λιγότερο σημεία τομής μπορούν να έχουν.\par
\possindeo Αν είναι όλες παράλληλες, $0$. Αν δύο είναι παράλληλες και η τρίτη τέμνει και τις δύο: $2$. Αν περνούν και οι τρεις από το ίδιο σημείο: $1$. Αν ανά δύο τέμνονται σε διαφορετικά σημεία: $3$.\par
\apantisi $0$, $1$, $2$ ή $3$ σημεία.\par
\elegxos Δεν μπορεί να υπάρχει περισσότερα από $\frac{3\cdot2}{2}=3$ σημεία (ένα ανά ζευγάρι).

\paradeigma{Παράλληλα τμήματα}
\tixero Οι ευθείες $\mathrm{ΑΒ}$ και $\Gamma\Delta$ είναι παράλληλες, με $\mathrm{ΑΒ}=3$ και $\Gamma\Delta=7$.\par
\tithelo Τι λέμε για τα τμήματα.\par
\possindeo Δύο τμήματα σε παράλληλες ευθείες είναι παράλληλα.\par
\apantisi $\mathrm{ΑΒ}\parallel\Gamma\Delta$, παρότι έχουν διαφορετικά μήκη.\par
\elegxos Η παραλληλία εξαρτάται από τις ευθείες και όχι από τα μήκη.

\tmima{Ασκήσεις}

\zoni{Το κατάλαβα}{παράλληλες και τεμνόμενες}

\ex Ποιες από τις παρακάτω ευθείες είναι παράλληλες και ποιες τεμνόμενες: α) οι ράγες του τρένου \quad β) τα δύο μονοπάτια του κήπου που διασταυρώνονται \quad γ) οι αράδες του τετραδίου \quad δ) δύο ευθείες που έχουν ένα κοινό σημείο.

\ex Σχεδίασε μια ευθεία $(\varepsilon)$ και ένα σημείο $\mathrm{Α}$ εκτός της. Φέρε από το $\mathrm{Α}$ την παράλληλη με γνώμονα και χάρακα.

\ex Πόσες παράλληλες προς μια ευθεία μπορούν να φέρουμε από ένα σημείο εκτός της; Από ένα σημείο πάνω της;

\ex Πόσα σημεία τομής μπορεί να έχουν $3$ διαφορετικές ευθείες; Σχεδίασε κάθε περίπτωση.

\ex Πόσα το πολύ σημεία τομής έχουν $4$ ευθείες; Και $\nu$ ευθείες;

\ex Αν $\varepsilon_1\parallel\varepsilon_2$ και $\varepsilon_2\parallel\varepsilon_3$, τι ξέρεις για τις $\varepsilon_1$, $\varepsilon_3$;

\ex Σωστό ή λάθος; α) Δύο ευθείες του επιπέδου που δεν τέμνονται είναι παράλληλες. \quad β) Δύο παράλληλα τμήματα έχουν ίσα μήκη. \quad γ) Αν δύο ευθείες έχουν δύο κοινά σημεία, ταυτίζονται. \quad δ) Από σημείο εκτός ευθείας περνούν δύο παράλληλες προς αυτήν.

\zoni{Το χρησιμοποιώ}{κήπος, ράγες, γήπεδο}

\ex \textbf{Οι σειρές.} Οι σειρές των λαχανικών είναι παράλληλες ευθείες και απέχουν $60$ cm. Πώς θα τις σημάνεις στο έδαφος με σχοινί και παλούκια;

\ex \textbf{Το γήπεδο.} Πόσα ζεύγη παράλληλων πλευρών έχει ένα ορθογώνιο γήπεδο; Πόσα ζεύγη κάθετων πλευρών;

\ex \textbf{Τα μονοπάτια.} Πέντε μονοπάτια του κήπου είναι ευθείες και δεν είναι δύο παράλληλα. Πόσα σημεία τομής έχουν το πολύ;

\ex \textbf{Οι ράγες.} Γιατί οι ράγες του τρένου πρέπει να είναι παράλληλες;

\ex \textbf{Το μοτίβο.} Σχεδίασε τρεις παράλληλες ευθείες και μια ευθεία που τις τέμνει. Πόσα σημεία τομής έχεις;

\zoni{Το σκέφτομαι}{ο Ευκλείδης και η παραλληλία}

\ex Γιατί από δύο σημεία περνά μία ευθεία ενώ από σημείο εκτός ευθείας περνά μία παράλληλη; Σε τι διαφέρουν οι δύο προτάσεις;

\ex Αν δύο ευθείες δεν έχουν κοινό σημείο στο χαρτί σου, μπορεί να τέμνονται σε σημείο πέρα από το χαρτί; Εξήγησε.

\ex Δύο διαφορετικές ευθείες είναι κάθετες σε μια τρίτη. Τι σχέση έχουν μεταξύ τους;

\ex Μπορείς να φέρεις δύο διαφορετικές παράλληλες από ένα σημείο προς μια ευθεία; Τι θα συνέβαινε;

\zoni{Για το σπίτι}{ψάξε παράλληλες}

\ex Βρες στο σπίτι πέντε ζεύγη παράλληλων ευθειών και τρία ζεύγη τεμνόμενων.

\ex Σχεδίασε με γνώμονα και χάρακα τρεις παράλληλες ευθείες και μέτρησε την απόσταση ανάμεσά τους.
